\documentclass[10pt,a4paper]{amsart}
\usepackage[margin=19mm]{geometry}
\usepackage{amsmath,amssymb,mathtools}
\usepackage[scr=rsfs,cal=euler]{mathalfa}
\usepackage{xfrac}
\usepackage[unicode,hidelinks,bookmarks=false]{hyperref}
\newcommand{\ie}{i.e.\ }
\newcommand{\cf}{cf.\ }
\newcommand{\resp}{respectively\ }
\newcommand{\Subsection}{\subsection}
\providecommand{\nobreakdash}{\nobreak\hbox{-}}
\setlength{\emergencystretch}{2em}
\multlinegap0pt
\usepackage{amssymb}
\usepackage{mathtools}
\usepackage{bm}
\usepackage[shortlabels]{enumitem}
\newtheorem{lemma}{Lemma}
\newtheorem{theorem}{Theorem}
\usepackage{cleveref}
\crefname{lemma}{Lemma}{Lemmas}
\crefname{theorem}{Theorem}{Theorems}
\newcommand{\Dplus}{\mathcal D_d^{+}}
\newcommand{\Htr}{\mathcal H_d^{0}}
\newcommand{\PhiK}{\Phi_K}
\newcommand{\Riem}{\mathrm{Riem}}
\DeclareMathOperator{\Tr}{Tr}
\DeclareMathOperator{\diag}{diag}
\title{A minimal qutrit counterexample to Conjecture~4.9 of Lesniewski and Ruskai}
\author{Domingos S. P. Salazar}
\date{Source: arXiv:2608.15464v1 (CC BY 4.0)}
\begin{document}
\maketitle

\noindent\emph{Source and licence.} A minimal qutrit counterexample to Conjecture~4.9 of Lesniewski and Ruskai. Version arXiv:2608.15464v1 (CC BY 4.0), distributed by arXiv under the Creative Commons Attribution 4.0 International licence, http://creativecommons.org/licenses/by/4.0/, as recorded in the arXiv licence metadata for this exact version and shown on the abstract page https://arxiv.org/abs/2608.15464v1. Full licence notice: \texttt{licenses/arxiv-2608.15464-CC-BY-4.0.txt}.\par\smallskip

\noindent\textit{Editorial scope.} Counted unit: the article's theorem thm:main (source0.tex lines 240--249). The article's complete original proof of it is delivered with it (source0.tex lines 251--303, one \texttt{proof} environment of 53 lines). No mathematics is written, repaired or re-derived: the statement and the proof are the article's own lines, in the article's own notation, and no retained line is substituted or re-flowed.  Retained uncounted as the source-local context the proof needs: lines 124--130; lines 133--139; lines 181--185; lines 209--212.  Nothing was omitted from any retained span, so the package carries no omission record.  No retained line contains a citation, so the wrapper carries no bibliography and no bibliography entry.  No retained line contains a figure environment, an image-inclusion command or drawing code, so no image is delivered and the package declares no asset dependency.  Editorial formatting only, in addition: an AMS \texttt{amsart} preamble that restates the article's own declarations and macros used by the retained lines, namely the environments \texttt{lemma}, \texttt{theorem} on separate counters, together with the standard packages those lines need.  The retained environments print in this document as Lemma 1, Theorem 1 in order of appearance, whereas the article prints the same items with the numbering of its own full document.  The complete native source of the article as distributed in the arXiv e-print for this exact version is bundled unchanged as \texttt{sources/207749/source0.tex}.
\begin{equation}\label{eq:eta}
 \eta_\kappa^{\Riem}(\Phi)
 =\sup_{\rho\in\Dplus}
  \sup_{\substack{A\in\Htr\\A\ne0}}
 \frac{\gamma_{\Phi(\rho)}^\kappa(\Phi(A),\Phi(A))}
      {\gamma_\rho^\kappa(A,A)},
\end{equation}

\begin{lemma}[Commuting reduction]\label{lem:commuting}
If $\rho\in\Dplus$ and $A\in\Htr$ commute, then for every $\kappa\in\mathcal K$,
\begin{equation}\label{eq:fisher}
 \gamma_\rho^\kappa(A,A)=\Tr \rho^{-1}A^2.
\end{equation}
Thus all normalized monotone metrics coincide with the Fisher--Rao metric on a commutative algebra.
\end{lemma}

\begin{equation}\label{eq:channel}
 \PhiK(X)
 =\sum_{i,j=1}^3
 K_{ij}\langle j|X|j\rangle\,|i\rangle\langle i|.
\end{equation}

\begin{equation}\label{eq:spectralvalue}
 \Lambda_2(\PhiK^\dagger\PhiK)
 =\frac{62+2\sqrt{61}}{225}.
\end{equation}

\begin{theorem}[Counterexample to Conjecture~4.9]\label{thm:main}
For the unital entanglement-breaking qutrit channel $\PhiK$ in \eqref{eq:channel} and every normalized monotone Riemannian metric $\kappa\in\mathcal K$,
\begin{equation}\label{eq:mainineq}
 \eta_\kappa^{\Riem}(\PhiK)
 \ge \frac{8896}{20007}
 >\frac{62+2\sqrt{61}}{225}
 =\Lambda_2(\PhiK^\dagger\PhiK).
\end{equation}
In particular, Lesniewski--Ruskai Conjecture~4.9 is false.
\end{theorem}

\begin{proof}
Choose the faithful diagonal state and traceless diagonal tangent
\begin{equation}\label{eq:witness}
 \rho_\star
 =\diag\!\left(\frac1{10},\frac45,\frac1{10}\right),
 \qquad
 A_\star=\diag(0,1,-1).
\end{equation}
Both the input pair and its image under $\PhiK$ commute. By \cref{lem:commuting}, the metric quotient is therefore independent of $\kappa$.

At the input,
\begin{equation}\label{eq:inputmetric}
 \gamma_{\rho_\star}^\kappa(A_\star,A_\star)
 =\frac{1}{4/5}+\frac{1}{1/10}
 =\frac{45}{4}.
\end{equation}
On diagonal entries, the channel acts as $K$, and
\begin{equation}\label{eq:images}
 K
 \begin{pmatrix}1/10\\[1mm]4/5\\[1mm]1/10\end{pmatrix}
 =\begin{pmatrix}1/10\\[1mm]19/50\\[1mm]13/25\end{pmatrix},
 \qquad
 K
 \begin{pmatrix}0\\[1mm]1\\[1mm]-1\end{pmatrix}
 =\begin{pmatrix}-2/3\\[1mm]2/5\\[1mm]4/15\end{pmatrix}.
\end{equation}
Thus
\begin{align}
 \gamma_{\PhiK(\rho_\star)}^\kappa
 \bigl(\PhiK(A_\star),\PhiK(A_\star)\bigr)
 &=\frac{(2/3)^2}{1/10}
   +\frac{(2/5)^2}{19/50}
   +\frac{(4/15)^2}{13/25} \notag\\
 &=\frac{11120}{2223}.
 \label{eq:outputmetric}
\end{align}
The corresponding admissible quotient in \eqref{eq:eta} is
\begin{equation}\label{eq:quotient}
 \frac{11120/2223}{45/4}
 =\frac{8896}{20007}.
\end{equation}
It remains only to verify the strict comparison with \eqref{eq:spectralvalue}. We use the rational intermediate value $26/75$:
\begin{equation}\label{eq:firstcompare}
 \frac{8896}{20007}-\frac{26}{75}
 =\frac{49006}{500175}>0,
\end{equation}
and
\begin{equation}\label{eq:secondcompare}
 \frac{26}{75}-\frac{62+2\sqrt{61}}{225}
 =\frac{16-2\sqrt{61}}{225}>0,
\end{equation}
because $\sqrt{61}<8$. Combining \eqref{eq:quotient}--\eqref{eq:secondcompare} proves \eqref{eq:mainineq} for every $\kappa\in\mathcal K$.
\end{proof}

\end{document}
